\documentclass[11pt]{amsart}
\usepackage{amsmath,amssymb}

\newtheorem{theorem}{Theorem}[section]
\newtheorem{lemma}[theorem]{Lemma}
\newtheorem{proposition}[theorem]{Proposition}
\newtheorem{corollary}[theorem]{Corollary}
\theoremstyle{definition}
\newtheorem{definition}[theorem]{Definition}

\title{A note on the nested radical iteration $x_{n+1}=\sqrt{2+x_n}$}
\author{Miriam K. Oduya}
\address{Department of Mathematics, Lakeshore State College, Halvorsen, MN, USA}
\email{m.oduya@lakeshore-state.example.edu}
\date{}

\begin{document}

\begin{abstract}
We give a self-contained treatment of the iteration $x_{n+1}=\sqrt{2+x_n}$
started at an arbitrary admissible point $a\ge -2$. We show that the iterates
always converge to $2$, with a geometric rate, and deduce the corresponding
statement for the Ces\`aro means.
\end{abstract}

\maketitle

\section{Introduction}

The nested radical $\sqrt{2+\sqrt{2+\sqrt{2+\cdots}}}$ is a standard first
example of a recursively defined sequence. In most textbooks the iteration is
started at $x_0=0$ or $x_0=\sqrt 2$, and convergence is obtained from the
monotone convergence theorem. The purpose of this note is to record the
argument for every starting value at which the iteration makes sense, that
is, for every $a\ge -2$, and to make the elementary estimates explicit.
Our main result is Theorem~\ref{thm:main}; in Section~\ref{sec:cesaro} we
derive a statement about averages from a classical theorem of Cauchy.

Throughout, $\mathbb N=\{0,1,2,\dots\}$, and $\sqrt{\cdot}$ denotes the
non-negative square root on $[0,\infty)$.

\section{The iteration}

\begin{definition}\label{def:seq}
Let $f\colon[-2,\infty)\to[0,\infty)$, $f(x)=\sqrt{2+x}$. For $a\ge -2$ we
define the sequence $(x_n)_{n\in\mathbb N}=(x_n(a))_{n\in\mathbb N}$ by
$x_0=a$ and $x_{n+1}=f(x_n)$ for $n\in\mathbb N$.
\end{definition}

Since $f$ takes values in $[0,\infty)\subset[-2,\infty)$, the sequence is
well defined, and $x_n\ge 0$ for all $n\ge 1$.

\begin{lemma}\label{lem:fixed}
For $x\ge -2$ we have $f(x)=x$ if and only if $x=2$.
\end{lemma}

\begin{proof}
Clearly $f(2)=\sqrt4=2$. Conversely, if $f(x)=x$ then $x\ge 0$, since $f$ is
non-negative, and squaring gives $x^2=x+2$, i.e.\ $(x-2)(x+1)=0$. As $x\ge 0$,
we get $x=2$.
\end{proof}

\begin{lemma}\label{lem:contr}
For every $x\ge -2$,
\[
  |f(x)-2|\le \tfrac14\,|x-2|.
\]
\end{lemma}

\begin{proof}
Rationalising,
\[
  f(x)-2=\frac{(2+x)-4}{\sqrt{2+x}+2}=\frac{x-2}{\sqrt{2+x}+2}.
\]
The denominator is at least $2$, whence the claim.
\end{proof}

\begin{lemma}\label{lem:mono}
If $-2\le x\le 2$, then $x\le f(x)\le 2$. If $x\ge 2$, then $2\le f(x)\le x$.
\end{lemma}

\begin{proof}
Let $-2\le x\le 2$. Then $2+x\le 4$, so $f(x)\le 2$. Since both sides are
non-negative, the inequality $x\le\sqrt{2+x}$ is equivalent to $x^2\le 2+x$,
that is, to $(x-2)(x+1)\le 0$, and this holds for $-2\le x\le 2$.

Now let $x\ge 2$. Then $2+x\ge 4$, so $f(x)\ge 2$; moreover $x\ge 2>0$, so
$f(x)\le x$ is equivalent to $2+x\le x^2$, i.e.\ to $(x-2)(x+1)\ge 0$, which
holds.
\end{proof}

\begin{corollary}\label{cor:monoseq}
Let $a\ge -2$ and $x_n=x_n(a)$.
\begin{enumerate}
\item If $a\le 2$, then $x_n\le x_{n+1}\le 2$ for all $n\in\mathbb N$.
\item If $a\ge 2$, then $2\le x_{n+1}\le x_n$ for all $n\in\mathbb N$.
\end{enumerate}
\end{corollary}

\begin{proof}
For (1), $x_0=a\in[-2,2]$, and if $x_n\in[-2,2]$ then Lemma~\ref{lem:mono}
gives $x_n\le x_{n+1}\le 2$; in particular $x_{n+1}\in[-2,2]$. Induction on
$n$ finishes the proof. Part (2) is identical, using the second half of
Lemma~\ref{lem:mono}.
\end{proof}

\begin{lemma}\label{lem:rate}
Let $a\ge -2$. Then $|x_n(a)-2|\le 2^{-n}|a-2|$ for every $n\in\mathbb N$.
\end{lemma}

\begin{proof}
For $n=0$ there is nothing to prove. If the inequality holds for $n$, then,
as $x_n\ge -2$, Lemma~\ref{lem:contr} gives
$|x_{n+1}-2|=|f(x_n)-2|\le\frac12|x_n-2|\le 2^{-(n+1)}|a-2|$.
\end{proof}

\begin{lemma}\label{lem:limcand}
Let $a\ge -2$ and suppose that $x_n(a)\to L$ for some $L\in\mathbb R$ with
$L\ne 2$. Then $L=-1$.
\end{lemma}

\begin{proof}
Since $x_n\ge -2$ for all $n$, $L\ge -2$, and $f$ is continuous on
$[-2,\infty)$. Letting $n\to\infty$ in $x_{n+1}=f(x_n)$ gives $L=\sqrt{2+L}$,
hence $L^2=L+2$ and $L\in\{2,-1\}$. As $L\ne 2$, $L=-1$.
\end{proof}

\begin{theorem}\label{thm:main}
For every $a\ge -2$ the sequence $(x_n(a))_{n\in\mathbb N}$ converges, and
\[
  \lim_{n\to\infty}x_n(a)=2.
\]
\end{theorem}

\begin{proof}
If $a\le 2$, then by Corollary~\ref{cor:monoseq}(1) the sequence is
non-decreasing and bounded above by $2$; if $a\ge 2$, it is non-increasing
and bounded below by $2$ by Corollary~\ref{cor:monoseq}(2). In either case the
monotone convergence theorem shows that $x_n\to L$ for some $L\ge -2$.
Passing to the limit in $x_{n+1}=f(x_n)$, using continuity of $f$ at $L$,
gives $f(L)=L$, so $L=2$ by Lemma~\ref{lem:fixed}. (Alternatively, if
$L\ne 2$ then $L=-1$ by Lemma~\ref{lem:limcand}, which is impossible because
$x_n\ge 0$ for $n\ge 1$ forces $L\ge 0$.) Lemma~\ref{lem:rate} gives the
same conclusion directly, together with a geometric rate.
\end{proof}

\section{Ces\`aro means}\label{sec:cesaro}

\begin{definition}\label{def:cesaro}
For a real sequence $y=(y_n)_{n\in\mathbb N}$ and $N\ge 1$, the $N$-th
Ces\`aro mean of $y$, i.e.\ the arithmetic mean of its first $N$ terms, is
\[
  \sigma_N(y)=\frac1N\sum_{n=0}^{N}y_n .
\]
\end{definition}

The following is the classical theorem of Cauchy on the regularity of the
Ces\`aro method; see~\cite{Cau21} and~\cite[Chapter~V]{Har49}.

\begin{theorem}[Cauchy]\label{thm:cauchy}
Let $(y_n)_{n\in\mathbb N}$ be a real sequence and $L\in\mathbb R$. Then
$y_n\to L$ if and only if $\sigma_N(y)\to L$ as $N\to\infty$.
\end{theorem}

\begin{corollary}\label{cor:cesaro}
Let $0\le a\le 2$ and $x_n=x_n(a)$. Then $\sigma_N(x)\to 2$ as
$N\to\infty$.
\end{corollary}

\begin{proof}
By Theorem~\ref{thm:main}, $x_n\to 2$. Theorem~\ref{thm:cauchy} now gives
$\sigma_N(x)\to 2$.
\end{proof}

\begin{thebibliography}{9}
\bibitem{Cau21} A.-L.~Cauchy, \emph{Cours d'analyse de l'\'Ecole royale
polytechnique. I\textsuperscript{re} partie: Analyse alg\'ebrique}, Debure,
Paris, 1821.
\bibitem{Har49} G.~H.~Hardy, \emph{Divergent Series}, Clarendon Press,
Oxford, 1949.
\end{thebibliography}

\end{document}
